\sectionkicker{Source-backed technical notes}
\subsection*{Agents \& Runtime — Codexを安全に長く動かす実行環境: Technical Notes}
本欄は記事本文で圧縮した一次資料上の情報を、比較・再検証しやすい形で整理したものである。時系列、確認済みの主張、留意点、一次資料URLを掲載する。
\medskip
\subsection*{Theme at a glance}
\addcontentsline{toc}{subsection}{Theme at a glance}
\begin{center}
\footnotesize
\begin{tabularx}{\linewidth}{@{}p{0.20\linewidth}p{0.10\linewidth}p{0.14\linewidth}X@{}}
\toprule
資料 & 位置づけ & 種別 & 時系列 \\
\midrule
Running Codex safely at OpenAI & 主要資料 & AGENT & 2026-05-08 (OFFICIAL\_PUBLICATION) \\
Building a safe, effective sandbox to enable Codex on Windows & 主要資料 & AGENT & 2026-05-13 (OFFICIAL\_PUBLICATION) \\
Kimi Code CLI May 2026 lifecycle & 主要資料 & OTHER & 2026-05 (PRODUCT\_RELEASE); 2026-05-26 (PRODUCT\_UPDATE); 2026-05-27 (PRODUCT\_UPDATE); 2026-05-28 (PRODUCT\_UPDATE); 2026-05-29 (PRODUCT\_UPDATE) \\
Mistral May 2026 technical releases & 主要資料 & OTHER & 2026-05-23 (AGENT\_RELEASE); 2026-05-22 (PRODUCT\_RELEASE); 2026-05-28 (AGENT\_RELEASE); 2026-05-27 (研究\_RELEASE) \\
Agentic AI Workload Characteristics & 補足資料 & 論文 & 2026-05-25 (論文\_RELEASE) \\
\bottomrule
\end{tabularx}
\end{center}
\normalsize

\begin{technicalnote}{Running Codex safely at OpenAI}{主要資料}
\begin{tabularx}{\linewidth}{@{}>{\bfseries}p{0.22\linewidth}X@{}}
組織 & OpenAI \\
種別 & AGENT \\
時系列 & 2026-05-08 (OFFICIAL\_PUBLICATION) \\
\end{tabularx}
\smallskip
{\bfseries 一次資料から整理したtechnical points}
\begin{itemize}[leftmargin=1.5em,itemsep=0.35em]
\item \textbf{一次情報で確認できる事実}: The preserved primary source identifies “Running Codex safely at OpenAI” on 2026-05-08.
\item \textbf{Vendor claim}: OpenAI's official feed describes its internal Codex safety setup using sandboxing, approvals, network policies, and agent-native telemetry.
\end{itemize}
{\bfseries 読む際の境界}
\begin{itemize}[leftmargin=1.5em,itemsep=0.35em]
\item \textbf{分析上の留意点}: The preserved first-party feed item is authoritative for publication chronology and vendor-described scope; detailed capability, benchmark, or safety claims were not independently reproduced.
\end{itemize}
{\bfseries 一次資料}
\begingroup\sloppy
\begin{itemize}[leftmargin=1.5em,itemsep=0.25em]
\item {\scriptsize\url{https://openai.com/index/running-codex-safely}}
\end{itemize}
\endgroup
\end{technicalnote}
\medskip

\begin{technicalnote}{Building a safe, effective sandbox to enable Codex on Windows}{主要資料}
\begin{tabularx}{\linewidth}{@{}>{\bfseries}p{0.22\linewidth}X@{}}
組織 & OpenAI \\
種別 & AGENT \\
時系列 & 2026-05-13 (OFFICIAL\_PUBLICATION) \\
\end{tabularx}
\smallskip
{\bfseries 一次資料から整理したtechnical points}
\begin{itemize}[leftmargin=1.5em,itemsep=0.35em]
\item \textbf{一次情報で確認できる事実}: The preserved primary source identifies “Building a safe, effective sandbox to enable Codex on Windows” on 2026-05-13.
\item \textbf{Vendor claim}: OpenAI's official feed describes a Windows Codex sandbox with controlled file access and network limits.
\end{itemize}
{\bfseries 読む際の境界}
\begin{itemize}[leftmargin=1.5em,itemsep=0.35em]
\item \textbf{分析上の留意点}: The preserved first-party feed item is authoritative for publication chronology and vendor-described scope; detailed capability, benchmark, or safety claims were not independently reproduced.
\end{itemize}
{\bfseries 一次資料}
\begingroup\sloppy
\begin{itemize}[leftmargin=1.5em,itemsep=0.25em]
\item {\scriptsize\url{https://openai.com/index/building-codex-windows-sandbox}}
\end{itemize}
\endgroup
\end{technicalnote}
\medskip

\begin{technicalnote}{Kimi Code CLI May 2026 lifecycle}{主要資料}
\begin{tabularx}{\linewidth}{@{}>{\bfseries}p{0.22\linewidth}X@{}}
組織 & Moonshot AI \\
種別 & OTHER \\
時系列 & 2026-05 (PRODUCT\_RELEASE); 2026-05-26 (PRODUCT\_UPDATE); 2026-05-27 (PRODUCT\_UPDATE); 2026-05-28 (PRODUCT\_UPDATE); 2026-05-29 (PRODUCT\_UPDATE) \\
\end{tabularx}
\smallskip
{\bfseries 一次資料から整理したtechnical points}
\begin{itemize}[leftmargin=1.5em,itemsep=0.35em]
\item \textbf{一次情報で確認できる事実}: The preserved primary source identifies “kimi-code-whats-new”.
\item \textbf{Vendor claim}: The Kimi Code changelog records its initial May release and rapid May 26–29 iteration, including OpenAI-compatible reasoning-model support, a plugin system, scheduled tasks, automatic permission handling, GitHub plugin installation, and removal of the default step limit.
\end{itemize}
{\bfseries 読む際の境界}
\begin{itemize}[leftmargin=1.5em,itemsep=0.35em]
\item \textbf{分析上の留意点}: The preserved first-party index is authoritative for the listed chronology, but capability/benchmark descriptions remain vendor claims and month-only dates retain month precision.
\end{itemize}
{\bfseries 一次資料}
\begingroup\sloppy
\begin{itemize}[leftmargin=1.5em,itemsep=0.25em]
\item {\scriptsize\url{https://www.kimi.com/code/docs/en/kimi-code/whats-new.html}}
\end{itemize}
\endgroup
\end{technicalnote}
\medskip

\begin{technicalnote}{Mistral May 2026 technical releases}{主要資料}
\begin{tabularx}{\linewidth}{@{}>{\bfseries}p{0.22\linewidth}X@{}}
組織 & Mistral AI \\
種別 & OTHER \\
時系列 & 2026-05-23 (AGENT\_RELEASE); 2026-05-22 (PRODUCT\_RELEASE); 2026-05-28 (AGENT\_RELEASE); 2026-05-27 (研究\_RELEASE) \\
\end{tabularx}
\smallskip
{\bfseries 一次資料から整理したtechnical points}
\begin{itemize}[leftmargin=1.5em,itemsep=0.35em]
\item \textbf{一次情報で確認できる事実}: The preserved primary source identifies “mistral-news”.
\item \textbf{Vendor claim}: The Mistral news snapshot records May technical releases including remote agents powered by Mistral Medium 3.5 on May 23, MCP-enabled Studio tooling and Workflows public preview on May 22, Vibe long-horizon Work/Code modes on May 28, and Physics AI announcements on May 27–28.
\end{itemize}
{\bfseries 読む際の境界}
\begin{itemize}[leftmargin=1.5em,itemsep=0.35em]
\item \textbf{分析上の留意点}: The preserved first-party index is authoritative for the listed chronology, but capability/benchmark descriptions remain vendor claims and month-only dates retain month precision.
\end{itemize}
{\bfseries 一次資料}
\begingroup\sloppy
\begin{itemize}[leftmargin=1.5em,itemsep=0.25em]
\item {\scriptsize\url{https://mistral.ai/news/}}
\end{itemize}
\endgroup
\end{technicalnote}
\medskip

\begin{technicalnote}{Agentic AI Workload Characteristics}{補足資料}
\begin{tabularx}{\linewidth}{@{}>{\bfseries}p{0.22\linewidth}X@{}}
組織 & - \\
種別 & 論文 \\
時系列 & 2026-05-25 (論文\_RELEASE) \\
\end{tabularx}
\smallskip
{\bfseries 一次資料から整理したtechnical points}
\begin{itemize}[leftmargin=1.5em,itemsep=0.35em]
\item \textbf{一次情報で確認できる事実}: The preserved primary source identifies “Agentic AI Workload Characteristics” on 2026-05-25.
\item \textbf{Author claim}: The paper characterizes ReAct-style agent workloads end-to-end, tracing repeated model calls, tools, and growing context across multiple benchmarks and model configurations.
\end{itemize}
{\bfseries 読む際の境界}
\begin{itemize}[leftmargin=1.5em,itemsep=0.35em]
\item \textbf{分析上の留意点}: The preserved original arXiv metadata/abstract is primary-paper evidence, but experimental results and generalization are author-reported and were not independently reproduced in this Evidence review.
\end{itemize}
{\bfseries 一次資料}
\begingroup\sloppy
\begin{itemize}[leftmargin=1.5em,itemsep=0.25em]
\item {\scriptsize\url{http://arxiv.org/abs/2605.26297v1}}
\end{itemize}
\endgroup
\end{technicalnote}
\medskip
